\section{Introduction}
\label{sec:introduction}

The Markoff equation\footnote{We use the spelling \emph{Markoff} for the
equation and its numbers, and refer to the uniqueness question as the
Markoff--Frobenius conjecture.}
$$
 a^2+b^2+c^2=3abc
$$
has a tree of positive integer solutions, generated from $(1,1,1)$ by
permutations and by the replacement $(a,b,c)\mapsto(a,b,3ab-c)$. The
Markoff--Frobenius conjecture asserts that the largest entry of a solution
determines the other two. It remains open, and a recurring strategy is to
attach to each occurrence of a largest entry in the tree an invariant finer
than the entry itself, rigid enough to separate two occurrences that share
their numerical value. The invariant studied here is a cost.

Each occurrence carries a reduced fraction $P/Q$, its Farey slope, and
determines a primitive representation of its label as a sum of two squares.
The subtractive Euclidean algorithm assigns an integer cost to such a pair:
the number of subtractions needed to bring the two coordinates to their
common divisor, or equivalently the sum of the partial quotients of their
ratio, minus one. For the occurrence pair that cost is exactly $P+Q-2$.
The natural question is whether the same integer admits a cheaper primitive
representation, and the comparison is subtle: a Gaussian reflection
preserves the norm while changing many partial quotients at once.

Reading the continued fraction of a source and applying the reflection is a
walk through a finite set of integer matrices, along which the cost
difference accumulates. Nonnegative weights on the cycles of such a graph
control repeated motion, but a finite walk also begins and ends somewhere,
and its initial and terminal charges survive any rebalancing of the edge
weights. Our first result concerns exactly that situation; the second
applies it to one Gaussian prime block.

\subsection{The two results}

The general theorem concerns a finite graph reading one binary letter on
each edge. Suppose an exact potential makes all initial, edge and terminal
weights nonnegative integers, and suppose each cyclic component of the
zero-weight subgraph carries a single letter frequency. Denote the finite
set of those frequencies by $\Theta$, assumed nonempty. For a balanced
accepted word of length $n$ with $m$ heavy letters and completed cost
$\Delta$, we prove
\begin{equation}
 \min_{\theta\in\Theta}|m-\theta n|<K(\Delta+1),
 \label{eq:intro-stability}
\end{equation}
where $K$ is obtained explicitly from the condensation graph and from
elementary count potentials inside its cyclic components. No
irreducibility, determinism, or bound on the input length is required.

The role of balance is precise. An unrestricted word can spend a long time
in one recurrent component and then a long time in another, so that its
overall frequency lies strictly between theirs. A balanced word cannot make
such a macroscopic mixture unless the two frequencies are already close.
This is what turns a statement about the convex hull of $\Theta$ into the
finite-set estimate \eqref{eq:intro-stability}.

Balance enters through a single compatible slope, valid simultaneously for
every factor. That slope has a geometric meaning: the pairs consisting of a
slope and an intercept that produce a given finite word form a convex region
of the plane, and the compatible slopes are exactly its shadow on the slope
axis. These regions are the faces of the arrangement through which Berstel
and Pocchiola counted the finite factors of Sturmian words, so the estimate
rests on one coordinate of a two-parameter picture.

Our arithmetic application fixes the Gaussian block $89=8^2+5^2$. If a
primitive source has norm divisible by $89$ exactly once, there is a single
complementary allocation of this pair of conjugate Gaussian primes, modulo
units and conjugation; we call the resulting comparison the complete
norm-$89$ flip. For the entire language $q=\varepsilon\chi(w)2$, where
$\chi$ doubles each digit, we prove that the flip cannot lower the cost,
and we determine every word on which the cost is preserved:
\begin{equation}
 q=111\chi(s)2,\qquad s\in(12111\mid21111)^*.
 \label{eq:intro-equality}
\end{equation}
The target exchanges the two five-letter tokens. The reduced-count
condition satisfied by a genuine Markoff occurrence excludes every nonempty
token sequence, so every distinct genuine target has strictly larger cost.

The same finite graph has critical frequencies
\begin{equation}
 \Theta_{89}=\{0,1/5,1/3,1\}.
 \label{eq:intro-critical}
\end{equation}
For a genuine source of Farey slope $P/Q$, the quantitative consequence is
\begin{equation}
 \min_{\theta\in\Theta_{89}}|P-\theta Q|\leqslant118\Delta+121.
 \label{eq:intro-farey}
\end{equation}
The excess therefore grows linearly with $Q$ on every closed interval of
slopes disjoint from $\Theta_{89}$. We construct an infinite sequence of
genuine sources to which this conclusion applies: their norms have $89$-adic
valuation exactly one, and both Farey coordinates tend to infinity.

\subsection{Background and contribution}

Finite transducers for homographic transformations of continued fractions
originate in Raney's work, and their matrix-state interpretation continues
to be useful in arithmetic applications
\cite{Raney1973,DongDupontODorneyWaitkus2025}. Potential transformations of
weighted automata, including the accompanying changes at initial and
terminal states, are standard \cite{Mohri2009}. Finiteness of an automaton
carries arithmetic information in other settings as well: over a finite
field, the power series that are algebraic are precisely those whose
coefficients a finite automaton can generate \cite{ChristolKamaeMendesFranceRauzy1980}.
We make no use of that equivalence here, and return to it in
\cref{sec:conclusion}. Critical classes and the
eventual concentration of optimal paths are central in min-plus spectral
theory \cite{Quadrat1997}, and the relation between cycle frequencies and
convex rotation sets has a well-developed graph-theoretic interpretation;
see \cite[Section 3]{BochiZhang2016}. We use these mechanisms explicitly and
attribute them where they occur. What \eqref{eq:intro-stability} adds is a
balance condition imposed on individual finite inputs, which replaces the
convex hull of the critical frequencies by the frequencies themselves.

The general statement proved here is the finite balanced-path estimate
\eqref{eq:intro-stability}, with its endpoint hypotheses and an explicit
condensation constant. We also record how narrow the arithmetic input is:
of the roughly $n^3/\pi^2$ balanced words of length $n$, at most $3n+3$ are
cores of genuine occurrences. The norm-$89$ application adds the endpoint
inequality, the complete equality language and the quantitative consequence
for genuine occurrences. We make no priority claim for the underlying
potential method, for the classical balance lemma, or for the
correspondence between central words and Markoff numbers.

The proof of the endpoint inequality at $89$ is in part computer-assisted:
it rests on finitely many integer inequalities attached to a finite graph,
verified in exact arithmetic. \Cref{sec:endpoint} states precisely which
finite assertion is verified and why it implies a statement about words of
arbitrary length. The stability theorem of \cref{sec:balanced-stability} and
the infinite family of \cref{sec:markoff-application} involve no computation
over words or labels.

The result compares a source with its complete norm-$89$ flip. A norm
containing other Gaussian prime blocks can have further primitive
representations, and no comparison with all of them is asserted here; nor
is the critical set \eqref{eq:intro-critical} claimed to be sharp among
genuine Markoff sources. These boundaries leave a concrete question for
further work: which recurrent zero-cost components can a central source
approach while retaining its arithmetic endpoint conditions?
